\subsection{Deep memory and compatibility of the cylinders}
\label{memprof:seccion}

Prolongation compares states that simultaneously retain residues,
quotients, and boundaries. In the six-coordinate linear chart of the
record, take the integer representative of \(L_0\) whose entries were
displayed above:
\[
 A_H=\begin{pmatrix}
 2&2&2&1&2&1\\2&1&2&2&1&1\\1&0&0&1&0&2\\
 0&0&1&1&1&0\\2&2&2&0&2&1\\2&1&2&1&0&0
 \end{pmatrix},\qquad\det A_H=1.
\]
The choice of this representative is part of the chart: reduction modulo
three retains \(L_0\), whereas the integer representative fixes the
operation on the deep quotients. With row vectors, this gives
\begin{equation}
 \begin{aligned}
 \mathsf{Div}_H:\mathbb Z^6&\longrightarrow
       \{0,1,2\}^6\times\mathbb Z^6,\\
 x_H&\longmapsto(u_H,x_H^+),\\
 y_H&=x_HA_H,\qquad u_H=[y_H]_3,\qquad
 x_H^+=(y_H-u_H)/3 .
 \end{aligned}
 \label{memprof:division}
\end{equation}
The pair \((u_H,x_H^+)\) retains \(y_H=u_H+3x_H^+\). Since
\(A_H\) is unimodular, its inverse has integer entries and
\begin{equation}
 x_H=(u_H+3x_H^+)A_H^{-1}.
 \label{memprof:inversa}
\end{equation}
This identity recovers the deep coordinate from its residue and quotient.
The other state fields retain their own records. The ternary orientation
is subsequently read through
\(0\mapsto0,\ 1\mapsto+1,\ 2\mapsto-1\), retaining the residue
\(u_H\in\{0,1,2\}^6\) and the quotient in \eqref{memprof:division}.

The cylinder chart relates a ternary prefix of length \(L\), with integer
index \(N\), to \(k_{\mathrm{dec}}\) decimal triads, with index \(D\).
The corresponding half-open cylinders are
\[
 I_3(N,L)=\left[\frac N{3^L},\frac{N+1}{3^L}\right),\qquad
 I_{1000}(D,k_{\mathrm{dec}})
 =\left[\frac D{1000^{k_{\mathrm{dec}}}},
         \frac{D+1}{1000^{k_{\mathrm{dec}}}}\right).
\]
Define the boundary integers
\[
 a_{\mathrm{cyl}}=N1000^{k_{\mathrm{dec}}}-D3^L,\qquad
 b_{\mathrm{cyl}}=(N+1)1000^{k_{\mathrm{dec}}}-D3^L.
\]
For a block of six residues, fix its integer encoding as
\[
 \nu(u_H)=\sum_{j=1}^6(u_H)_j3^{6-j}
 \in\{0,\ldots,728\}.
\]

\begin{lemma}[Integer update of compatible cylinders]
\label{memprof:cilindros}
The preceding cylinders intersect if and only if
\(a_{\mathrm{cyl}}<3^L\) and \(b_{\mathrm{cyl}}>0\).
Appending a ternary block with index
\(\nu\in\{0,\ldots,728\}\) gives \(N'=729N+\nu\) and
\(L'=L+6\). For \(\Delta k_{\mathrm{dec}}\in\{0,1\}\),
the boundaries update according to
\[
 \begin{array}{ll}
 \Delta k_{\mathrm{dec}}=0:
   &a_{\mathrm{cyl}}'=729a_{\mathrm{cyl}}+\nu1000^{k_{\mathrm{dec}}},\\[2pt]
 \Delta k_{\mathrm{dec}}=1:
   &D'=1000D+\delta,\\
   &a_{\mathrm{cyl}}'=729000a_{\mathrm{cyl}}
      +\nu1000^{k_{\mathrm{dec}}+1}-\delta\,729\,3^L ,
 \end{array}
\]
where \(\delta\in\{0,\ldots,999\}\) is the triad of the compatible
prolongation. In both cases,
\[
 b_{\mathrm{cyl}}'=a_{\mathrm{cyl}}'
          +1000^{k_{\mathrm{dec}}+\Delta k_{\mathrm{dec}}}.
\]
\end{lemma}
\begin{proof}
Two half-open intervals intersect precisely when each left endpoint is
smaller than the right endpoint of the other interval. Multiplication
of these two inequalities by \(3^L1000^{k_{\mathrm{dec}}}\) gives
\(a_{\mathrm{cyl}}<3^L\) and \(b_{\mathrm{cyl}}>0\).
Substitution of \(N'\), \(L'\), and \(D'\) into
\[
 a_{\mathrm{cyl}}'
 =N'1000^{k_{\mathrm{dec}}+\Delta k_{\mathrm{dec}}}-D'3^{L'}
\]
gives the two recurrences. Subtracting the definitions of
\(b_{\mathrm{cyl}}'\) and \(a_{\mathrm{cyl}}'\) yields the final identity.
\end{proof}

These operations act on integer prefixes and their compatible
restrictions. The Archimedean reading appears when the resulting system
of cylinders is evaluated. The prolongation domain also incorporates
orientation, carry, the joint signature, and the boundary data of the state
\(\widehat X\) in \eqref{eq:actualizacion}; interval comparison is one
component of this joint compatibility.
